\section{Coefficient transitions preserve the original relative geometry}

\begin{lemma}[An arbitrary arithmetic coefficient extension]
\label{lem:kn-relative-arithmetic-scalar-transition}
\lean{EnvelopeProperty.Vanishing.RelativeArithmeticScalarTransition.relative_isPullback,EnvelopeProperty.Vanishing.RelativeArithmeticScalarTransition.isAmple,EnvelopeProperty.Vanishing.RelativeArithmeticScalarTransition.original_isPullback}
\leanok
Let $g:Y\to\operatorname{Spec}D$ have composite
$f:Y\to\operatorname{Spec}A$, and let $S$ be any $A$-algebra.
The original coefficient extension
$Y_S=Y\times_A S$ has a canonical map
$g_S:Y_S\to\operatorname{Spec}(S\otimes_A D)$ which is the base
change of $g$.  Its composite to $\operatorname{Spec}S$ is the
original second projection.  Every original relatively ample line
pulls back to a relatively ample line for $g_S$.  If an actual
cartesian original-field square $X=Y\times_A k$ is given and
$S\subset k$ contains $A$, the canonical original map
$X\to Y_S$ has the original field structure and is cartesian;
successive module pullbacks recover the original pullback from $Y$.
\end{lemma}
\begin{proof}
The tensor affine square identifies
$\operatorname{Spec}(S\otimes_A D)$ with
$\operatorname{Spec}S\times_A\operatorname{Spec}D$.
The two original projections of $Y_S$, with the first followed by
$g$, induce $g_S$.  Its composites are the original projections by
the product universal property.  Pasting with the tensor affine
square gives the defining square for $Y_S$; cancellation shows
that $g_S$ is the base change of $g$.  Pull back the actual closed
projective embedding and original positive power witnessing
relative ampleness.  Projective space base change and compatibility
of module tensor powers with pullback give the same positive-power
criterion on $Y_S$.

For the original field, the map $X\to Y_S$ is induced by
$X\to Y$ and $X\to\operatorname{Spec}k\to\operatorname{Spec}S$.
Their coefficient composites agree by the original cartesian
square.  Pasting its proposed field square with the defining square
of $Y_S$ gives that original cartesian square.  Cancellation again
proves the new field square cartesian.  The first projection of
$X\to Y_S$ is the original map to $Y$, so associativity of module
pullback gives the stated whole-module recovery isomorphism.
\end{proof}
